\documentclass[12pt,a4paper]{article}

\usepackage[a4paper,left=1in,right=1in,top=1in,bottom=1in]{geometry}
\usepackage[T1]{fontenc}
\usepackage{newtxtext}
\usepackage{newtxmath}
\usepackage[bahasa]{babel}
\usepackage{setspace}
\usepackage{indentfirst}
\usepackage{booktabs}
\usepackage{longtable}
\usepackage{graphicx}
\usepackage{float}
\usepackage[colorlinks=true,linkcolor=blue,urlcolor=blue]{hyperref}
\usepackage{caption}
\captionsetup{font=small,labelsep=period}

\setstretch{1.5}

\title{\textbf{Laporan Proyek Expert Advisor (EA)\\[0.5em]Moving Average Crossover untuk MetaTrader 5}}
\author{Rayhan Haldi Hermawan\\ NIM 24/545406/PA/23176}
\date{}

\begin{document}

\maketitle
\thispagestyle{empty}
\newpage
\pagenumbering{arabic}

\begin{center}
\section*{Laporan Proyek Sains Manajemen 261}
\vspace{0.5em}
{\normalsize Moving Average Crossover Expert Advisor (EA) untuk MetaTrader 5}
\end{center}

\section{Pendahuluan}

\subsection{Latar Belakang}
Perdagangan otomatis (automated trading) semakin banyak digunakan oleh para trader untuk
menghilangkan faktor emosi dan meningkatkan konsistensi eksekusi strategi. Salah satu
platform yang mendukung perdagangan otomatis adalah MetaTrader 5 (MT5) dengan bahasa
pemrograman MQL5. Dalam proyek ini, sebuah Expert Advisor (EA) berbasis strategi
\textit{Moving Average Crossover} dibuat, dikompilasi, dan siap diuji menggunakan
\textit{Strategy Tester} MT5.

\subsection{Tujuan}
Tujuan dari proyek ini adalah:
\begin{enumerate}
\item Membuat Expert Advisor sederhana untuk MT5 menggunakan bahasa MQL5.
\item Menerapkan strategi \textit{Moving Average Crossover} sebagai dasar logika sinyal.
\item Melakukan uji coba (\textit{backtest}) dan optimasi parameter pada \textit{Strategy Tester}.
\item Membuat laporan hasil uji coba sebagai dokumentasi proyek.
\end{enumerate}

\section{Tinjauan Referensi}
Strategi pada proyek ini mengacu pada kanal YouTube
\textit{Ren\'e Balke (BM Trading)}, khususnya video tutorial
\textit{"Fully working Moving Average MT5 Expert Advisor Programming Tutorial"}
yang menjelaskan pembuatan EA bergerak rata-rata pada MT5.

\begin{itemize}
\item Nama Channel: Ren\'e Balke (BM Trading)
\item Link Video: \url{https://youtu.be/T78Q7K3c11s}
\item Link Kode (MQ5): \url{https://github.com/VelAstra/Sains-Manajemen-10-EA-dengan-AI/blob/main/eas/MovingAverageCrossover.mq5}
\end{itemize}

\section{Metodologi}

\subsection{Konsep Strategi Moving Average Crossover}
Moving Average (MA) adalah indikator yang menghitung rata-rata harga dalam periode tertentu
untuk memperhalus pergerakan harga. Strategi \textit{crossover} menggunakan dua MA dengan
periode berbeda, yaitu MA cepat dan MA lambat. Sinyal dihasilkan saat kedua MA saling
berpotongan:

\begin{description}
\item[Cross Up (Buy)] terjadi ketika MA cepat memotong MA lambat dari bawah ke atas,
\item[Cross Down (Sell)] terjadi ketika MA cepat memotong MA lambat dari atas ke bawah.
\end{description}

\subsection{Pengaturan Parameter}
Parameter EA yang dapat diatur pada proyek ini adalah sebagai berikut:

\begin{longtable}{@{}p{0.42\textwidth}p{0.2\textwidth}p{0.3\textwidth}@{}}
\toprule
\textbf{Parameter} & \textbf{Nilai Default} & \textbf{Keterangan} \\
\midrule
Periode MA cepat & 10 & Periode moving average cepat \\
Periode MA lambat & 50 & Periode moving average lambat \\
Metode MA & SMA & SMA / EMA / SMMA / LWMA \\
Harga acuan & Close & Harga yang digunakan indikator \\
Ukuran lot & 0.01 & Volume transaksi \\
Stop loss & 0 (nonaktif) & Jarak SL dalam poin \\
Take profit & 0 (nonaktif) & Jarak TP dalam poin \\
Batas spread maksimum & 50 & Poin maksimal spread \\
Magic number & 1001 & Penanda posisi milik EA \\
\bottomrule
\end{longtable}

\subsection{Keamanan Akun}
EA dilengkapi pengaman agar tidak berjalan pada akun \textit{real}. Jika dipasang pada
akun real dan parameter \texttt{InpAllowReal} bernilai \texttt{false}, maka EA menolak
untuk berjalan. Hal ini dilakukan untuk melindungi dana dari risiko yang tidak diinginkan
selama masa pengujian.

\section{Implementasi}

\subsection{Fitur Utama EA}
\begin{enumerate}
\item Mendeteksi sinyal hanya pada bar baru agar tidak terjadi sinyal ganda.
\item Membuka maksimal satu posisi dalam satu waktu (\textit{one position at a time}).
\item Menolak membuka posisi jika spread melebihi batas yang ditentukan.
\item Menyetel stop loss dan take profit secara otomatis jika diaktifkan.
\item Menggunakan \texttt{magic number} untuk membedakan posisi milik EA.
\end{enumerate}

\subsection{Pseudokode}
\begin{verbatim}
Awali EA (OnInit):
    Periksa jenis akun (real/demo)
    Buat indikator MA cepat dan MA lambat

Setiap tick (OnTick):
    Jika belum ada bar baru, keluar
    Jika spread terlalu besar, keluar
    Ambil nilai MA cepat & lambat (bar sekarang dan bar sebelumnya)

    Cross up   = MA_cpt(prev) <= MA_lmb(prev) && MA_cpt(curr) > MA_lmb(curr)
    Cross down = MA_cpt(prev) >= MA_lmb(prev) && MA_cpt(curr) < MA_lmb(curr)

    Jika sudah ada posisi, keluar
    Jika cross up,  buka posisi BUY
    Jika cross down, buka posisi SELL
    Terapkan SL/TP jika diaktifkan
\end{verbatim}

\section{Hasil Uji Coba (Backtest)}

\subsection{Pengaturan Pengujian}
Pengujian dilakukan menggunakan \textit{Strategy Tester} MetaTrader 5 sebagai berikut:

\begin{longtable}{@{}p{0.42\textwidth}p{0.5\textwidth}@{}}
\toprule
\textbf{Aspek} & \textbf{Setting} \\
\midrule
Expert Advisor & MovingAverageCrossover \\
Simbol & [ISI: contoh EURUSD] \\
Timeframe & [ISI: contoh M15] \\
Model & Every tick \\
Modal awal & [ISI] \\
Periode pengujian & [ISI: contoh 2024-01-01 s.d. 2024-06-30] \\
\bottomrule
\end{longtable}

\subsection{Hasil Pengujian Awal}
Pada [ISI], hasil backtest awal (parameter default MA 10/50) adalah sebagai berikut:

\begin{longtable}{@{}p{0.6\textwidth}p{0.3\textwidth}@{}}
\toprule
\textbf{Metrik} & \textbf{Hasil} \\
\midrule
Total net profit & [ISI] \\
Profit factor & [ISI] \\
Total trades & [ISI] \\
Jumlah buy / sell & [ISI] / [ISI] \\
Persentase profitable & [ISI] \\
Maximum drawdown & [ISI] \\
Angka Sharpe \textit{(jika tersedia)} & [ISI] \\
\bottomrule
\end{longtable}

\newpage
\subsection{Optimasi Parameter}
Optimasi dilakukan untuk mencari kombinasi parameter terbaik menggunakan fitur
\textit{Strategy Optimization} bawaan MT5. Parameter yang menjadi kandidat optimasi antara
lain periode MA cepat, periode MA lambat, metode MA, serta stop loss dan take profit.

\begin{longtable}{@{}p{0.5\textwidth}p{0.2\textwidth}p{0.2\textwidth}@{}}
\toprule
\textbf{Area Optimasi} & \textbf{Rentang} & \textbf{Langkah} \\
\midrule
Periode MA cepat & [ISI] & [ISI] \\
Periode MA lambat & [ISI] & [ISI] \\
Metode MA & [ISI] & --- \\
Stop loss (poin) & [ISI] & [ISI] \\
Take profit (poin) & [ISI] & [ISI] \\
\bottomrule
\end{longtable}

Hasil kombinasi terbaik (misalnya berdasarkan \textit{profit factor} atau
\textit{expected payoff}):

\begin{longtable}{@{}p{0.6\textwidth}p{0.3\textwidth}@{}}
\toprule
\textbf{Parameter} & \textbf{Nilai Terbaik} \\
\midrule
Periode MA cepat & [ISI] \\
Periode MA lambat & [ISI] \\
Metode MA & [ISI] \\
Stop loss & [ISI] \\
Take profit & [ISI] \\
Profit factor (hasil) & [ISI] \\
\bottomrule
\end{longtable}

\section{Kesimpulan}
[ISI KESIMPULAN: misalnya proyek berhasil membuat EA Moving Average Crossover pada MT5;
EA berhasil dikompilasi tanpa error dan siap diuji pada Strategy Tester dengan uang
virtual (simulasi). Hasil backtest dan optimasi menunjukkan ...]

\section{Referensi}
\begin{enumerate}
\item Balke, Ren\'e. \textit{BM Trading --- Free Expert Advisors for MetaTrader 5.}
      Diambil dari \url{https://en.bmtrading.de} (diakses [ISI]).
\item Ren\'e Balke. \textit{Fully working Moving Average MT5 Expert Advisor Programming
      Tutorial.} YouTube. \url{https://youtu.be/T78Q7K3c11s} (diakses [ISI]).
\item MetaQuotes. \textit{MQL5 Reference.} \url{https://www.mql5.com/en/docs}
      (diakses [ISI]).
\end{enumerate}

\end{document}